\section{Design Decisions}
\label{sec:decisions}

Each subsection follows the same shape: the options on the table, the choice,
and the measurement that forced it.

\subsection{D1: The testbed --- why not a popular open-source chip}
\label{sec:d1}

\textbf{Options.} (a)~Use OpenTitan, the best-documented open SoC, as-is;
(b)~apply a consistent renaming/constant-perturbation transform to OpenTitan
to synthesize a chip that provably exists nowhere (``Cisco-X''); (c)~use a
real but obscure platform.

\textbf{Choice.} (c), the PULP Carfield/Cheshire platform, with (b) validated
and kept as a companion methodology.

\textbf{Why.} A knowledge-injection demo is meaningless if the model already
knows the answers. We probed Claude Opus~5 closed-book on 252 auto-generated
OpenTitan questions: 75.0\% correct, including exact register offsets---the
chip is \emph{in the weights}, so option~(a) cannot attribute anything to
training. The Cisco-X transform (identifier renaming through the register-tool
regeneration flow, constant perturbation, layout shuffling) drives frontier
closed-book performance to 0.0\% with 99.6\% explicit abstention while
open-book performance stays at 91.8\%---clean attribution, and a reusable
trick for customer engagements where even question texts are sensitive. But a
synthetic chip cannot demonstrate handling of \emph{real} corpus mess:
Carfield is actively developed, sparsely discussed on the public web
(Opus~5: 72.0\% on our final benchmark, with register offsets at 18\%), and
comes with the full organic sprawl---259 repositories, generated headers,
issues, drivers---that a customer's internal data will have.

\subsection{D2: Adaptation method --- DAPT before SFT, weights before retrieval}
\label{sec:d2}

\textbf{Options.} (a)~Retrieval over the corpus with a strong general model;
(b)~supervised fine-tuning on Q/A pairs; (c)~domain-adaptive continued
pretraining (DAPT), optionally followed by SFT.

\textbf{Choice.} (c), DAPT first; SFT deferred until knowledge is in the
weights.

\textbf{Why.} The property under test is closed-book knowledge
(\S\ref{sec:problem}); retrieval by construction does not put knowledge in
weights, so (a) answers a different question. SFT-first (b) teaches answer
\emph{format} but must invent its Q/A pairs from the same corpus---and our
corpus is 31.2M tokens of mostly non-Q/A material; at that scale the
ChipNeMo-style recipe~\citep{liu2023chipnemo} (full-parameter continued
pretraining at low learning rate with a general-data replay mix) is the
established path, and our token budget sits comfortably in the regime where
continued pretraining is worthwhile. The 3-shot completion protocol
substitutes for instruction-following during evaluation, letting us measure
injected knowledge without an SFT stage confounding it.

\subsection{D3: Base model --- a base checkpoint, not a chat model}
\label{sec:d3}

\textbf{Options.} (a)~Instruction-tuned 7--9B checkpoints; (b)~base
checkpoints.

\textbf{Choice.} Qwen3.5-9B-\emph{Base}, full-parameter, LR $5\times10^{-6}$
cosine, packing at 4096, two epochs.

\textbf{Why.} ChipNeMo's ablations are unambiguous---DAPT on a chat model
degrades both knowledge uptake and chat ability; low learning rate and
full-parameter updates (not LoRA) are what make fact absorption work---and we
saw no reason to relitigate them on our budget. We took the newest open base
model at the largest size our 4-GPU budget trains in about an hour.

\subsection{D4: The data recipe --- where the result actually comes from}
\label{sec:d4}

This is the decision that produced the paper's central lesson, in three
measured steps.

\textbf{Step 1: raw corpus fails.} The obvious recipe---feed the filtered
31.2M-token organizational corpus (RTL, generated headers, docs, drivers,
issues of all 259 non-fork repositories, plus 8\% general-text replay) to the
trainer---\emph{looks} like it works: training loss falls from 0.80 to 0.35.
It injects nothing extractable: on the pilot benchmark's 42 register-offset
questions the raw-DAPT model scores 1/42, identical to base. A fact that
appears in exactly one surface form (a \texttt{\#define} line) is stored but
not retrievable as an answer, exactly as \citet{allenzhu2023physics31} predict
with synthetic biographies; we believe this is the clearest industrial-corpus
confirmation of that result to date, and it means the popular mental model
``CPT on your data $\Rightarrow$ your model knows your data'' fails
\emph{silently}---loss curves report success while extraction stays at zero.

\textbf{Step 2: augmentation works exactly as far as it covers.} We then
restated facts in multiple declarative paraphrases mixed into the corpus. The
first version covered only the 461 facts our pilot questions quizzed---and its
gains land \emph{only} on layers those facts populate. On the final benchmark,
that model (55.2\%) matches the untrained base \emph{exactly} on everything
its augmentation never covered: dependency pins 0/8, region sizes 1/8,
driver--register links 4/15. Coverage determines the extractable scope;
evaluation must therefore be a random audit of the full fact inventory, never
the augmentation's target list.

\textbf{Step 3: diversity, reversal, and context against interference.} The
final generator enumerates \emph{every} fact from every structured source
($\sim$340 register offsets, all register-description fields, both memory
maps, build pins, driver--register links from static analysis, issue titles)
and emits each through 24 LLM-written templates per fact type (generated for
\$0.30 total, validated programmatically), of which at least a third are
\emph{reversed}---value first, entity second---to dodge the reversal
curse~\citep{berglund2023reversal}. Two further rules come from an
interference observation: at identical exposure, 9 mutually distinctive base
addresses were memorized 8/9 while 339 near-identical offsets managed 9/42.
Similarity, not exposure count, was the bottleneck. The countermeasures are
whole-table narrative documents (the register map rendered as ascending prose,
a markdown table, a C header, and a neighbor-chain walk, so that similar facts
appear as each other's context) and fewer repeats (6 rather than 8) in favor
of more forms. The augmentation totals 3.49M tokens, about 10\% of the
training mix. Result: register-offset recall went from 1/42 (raw) and 9/42
(quizzed-facts augmentation) to 25/42 on the pilot protocol---and 97\%
(323/333) over \emph{all} offsets under the final protocol.

\subsection{D5: Evaluation --- auditing our own benchmark}
\label{sec:d5}

\textbf{Options.} (a)~Quiz whatever is easy to auto-generate; (b)~design the
distribution top-down from what an engineer needs to know, and treat the
generator itself as a system under test.

\textbf{Choice.} (b), after (a) burned us.

\textbf{Why.} Our pilot benchmark drew 54\% of its questions from register
offsets---not because offsets matter most, but because a regex over generated
headers made them free. Auto-generated evaluations inherit the streetlight
effect of their sources, and training data designed against such an
evaluation inherits it twice. The final benchmark is layered by knowledge
kind: L1 facts (30\%), L2 structure---memory-map ownership, sizes, dependency
pins (20\%), L3 behavior---documentation semantics as 4-way multiple choice
(25\%), L4 cross-source---which driver function touches which register, from
static analysis (15\%), L5 engineering history (10\%); every question is
machine-checkable.

Three generator defects surfaced only because we ran a frontier model as an
adversarial probe, and each became a construction rule. \emph{Lexical leak:}
Opus scored 100\% on description$\rightarrow$register-name multiple choice by
word overlap alone (``Ethernet RGMII PHY clock divider''
$\rightarrow$ \texttt{ETH\_RGMII\_PHY\_CLK\_DIV}); we now keep a
multiple-choice question only if every distractor's lexical-overlap score is
at least the answer's, and we deleted one question type (RTL
instance$\rightarrow$module) that could not be repaired. \emph{Few-shot
format holes:} a shot set whose three exemplars were all bit-range answers
caused models to answer hex questions in decimal---mis-scored by the
normalizer \emph{and} genuinely degraded---silently depressing an entire
protocol; every answer format now has an exemplar. \emph{Worthless question
types:} issue-number$\rightarrow$title recall is a pure arbitrary mapping
with no engineering value; no model exceeded 27\%, and we removed the type
(1{,}423 questions) from the bank rather than let it dominate the aggregate.

\subsection{D6: Systems configuration --- the cheap 2.2$\times$}
\label{sec:d6}

Three configuration-level choices, found by 15-step smoke tests, cut cost
without touching the science: the \texttt{flash-linear-attention} kernels are
mandatory for Qwen3.5's gated-delta-net layers (25\,s/it $\rightarrow$
7.4\,s/it); trading ZeRO-3's per-layer parameter gathering for either ZeRO-2
or a micro-batch of 2~\citep{rajbhandari2019zero} gives an equivalent further
2.2$\times$ (7.43 $\rightarrow$ 3.41\,s/it), making the final training run 77
minutes on 4$\times$H100 ($\approx$\$20); and evaluation under
vLLM~\citep{kwon2023vllm} with prefix caching amortizes the shared few-shot
prompt across the whole bank. Training used
LLaMA-Factory~\citep{zheng2024llamafactory} with DeepSpeed; the allocator
setting \texttt{expandable\_segments} suppressed fragmentation warnings that
otherwise threaten long runs.
